\subsection{\texorpdfstring{Compact convergence in \(\mathcal{M}(\mathrm{X};\mathbf{C})\)}{Compact convergence in M(X; C)}}

Recall that the topology of compact convergence on \(\mathcal{M}(\mathrm{X};\mathbf{C})\) is the topology of uniform convergence in the compact subsets of \(\mathcal{K}(\mathrm{X};\mathbf{C})\) . We shall call topology of strictly compact convergence on \(\mathcal{M}(\mathrm{X};\mathbf{C})\) the topology of uniform convergence in the strictly compact subsets (No.~1) of \(\mathcal{K}(\mathrm{X};\mathbf{C})\) .

PROPOSITION 17. --- On the space \(\mathcal{M}(\mathrm{X};\mathbf{C})\) , consider the following topologies:

\(\mathcal{T}_1\) : the topology of pointwise convergence in a total subset \(\mathrm{T}\) of \(\mathcal{K}(\mathrm{X};\mathbf{C})\) ;

\(\mathcal{T}_2\) : the vague topology;

\(\mathcal{T}_3\) : the topology of strictly compact convergence;

\(\mathcal{T}_4\) : the topology of compact convergence.

Each of these topologies is coarser than the next. Moreover:

\begin{enumerate}
\def\labelenumi{(\roman{enumi})}
\item
  The bounded sets are the same for \(\mathcal{T}_2\) , \(\mathcal{T}_3\) and \(\mathcal{T}_4\) .
\item
  If \(\mathbf{H}\) is a vaguely bounded subset of \(\mathcal{M}(\mathrm{X};\mathbf{C})\) , the topologies induced on \(\mathbf{H}\) by the topologies \(\mathcal{T}_1, \mathcal{T}_2, \mathcal{T}_3, \mathcal{T}_4\) are identical.
\end{enumerate}

A vaguely bounded subset H of \(\mathcal{M}(X;C)\) is equicontinuous (No.~9, Prop. 15), thus the first assertion follows from TVS, III, §3, No.~7, Prop. 9, and the second from GT, X, §2, No.~4, Th. 1.

Recall that when X is paracompact, the topology of strictly compact convergence coincides with the topology of compact convergence (No.~1, Prop. 2).

PROPOSITION 18. --- On the cone \(\mathcal{M}_{+}(X)\) , the topologies induced by the following topologies coincide:

\(\mathcal{T}_{1}\) : the topology of pointwise convergence in a linear subspace V of \(\mathcal{K}(X;C)\) that is dense in \(\mathcal{K}(X;C)\) and satisfies the property (P) (No.~7, Prop. 9);

\(\mathcal{T}_2\) : the vague topology;

\(\mathcal{T}_3\) : the topology of strictly compact convergence.

Since every filter is the intersection of the ultrafilters finer than it (GT, I, §6, No.~4, Prop. 7), it suffices to show that if \(\mathfrak{U}\) is an ultrafilter on \(\mathcal{M}_{+}(\mathrm{X})\) that converges to a measure \(\mu_0\) for the topology \(\mathcal{T}_1\) , then it also converges to \(\mu_0\) for \(\mathcal{T}_3\) . Let \(\mathbf{K}\) be a compact subset of \(\mathbf{X}\) ; by hypothesis, there exists a function \(h \in \mathrm{V}\) that is \(\geqslant 0\) on \(\mathbf{X}\) and takes values \(>0\) on \(\mathbf{K}\) ; it follows that every function \(f \in \mathcal{K}(\mathrm{X},\mathrm{K};\mathbf{C})\) may be written \(f = gh\) with \(g \in \mathcal{K}(\mathrm{X},\mathrm{K};\mathbf{C})\) , and if \(c = \inf_{x \in \mathrm{K}} h(x) > 0\) then \(\| g \| \leqslant c^{-1} \| f \|\) . By hypothesis, there exists a set \(\mathrm{H}_0 \in \mathfrak{U}\) such that, for every measure \(\mu \in \mathrm{H}_0\) ,

\[
0 \leqslant \mu (h) \leqslant \mu_ {0} (h) + 1 = b.
\]

Consequently, for every function \(f \in \mathcal{K}(\mathrm{X};\mathbf{C})\) ,

\[
| \langle f, h \cdot \mu \rangle | = | \langle h f, \mu \rangle | \leqslant \| f \| \cdot \mu (h) \leqslant b \| f \|
\]

for every measure \(\mu\in\mathrm{H}_{0}\) ; this proves that the set H of measures \(h\cdot\mu\) , where \(\mu\) runs over \(\mathrm{H}_{0}\) , is vaguely bounded. If \(\mathfrak{U}_{0}\) is the ultrafilter induced by \(\mathfrak{U}\) on \(\mathrm{H}_{0}\) , the image of \(\mathfrak{U}_{0}\) under the mapping \(\mu\mapsto h\cdot\mu\) is the base of an ultrafilter \(\mathfrak{F}\) on H, and since H is relatively compact for the topology of strictly compact convergence (Prop. 17 and No.~9, Prop. 15), \(\mathfrak{F}\) is convergent to a measure \(\nu_{0}\) for this topology. In other words, for any \(\varepsilon>0\) and any compact subset L of \(\mathcal{K}(\mathrm{X},\mathrm{K};\mathbf{C})\) , there exists a subset N of \(\mathrm{H}_{0}\) belonging to \(\mathfrak{U}\) such that, for every function \(g\in\mathrm{L}\) and every pair of measures \(\mu,\mu'\) belonging to N, one has \(|\langle g,h\cdot\mu\rangle-\langle g,h\cdot\mu'\rangle|\leqslant\varepsilon\) , that is,

\[
| \langle g h, \mu \rangle - \langle g h, \mu^ {\prime} \rangle | \leqslant \varepsilon .
\]

Now, we saw above that the mapping \(g \mapsto gh\) is an automorphism of the Banach space \(\mathcal{K}(\mathrm{X},\mathrm{K};\mathbf{C})\) . We have thus shown that U is a Cauchy filter on \(\mathcal{M}_{+}(\mathrm{X})\) for the topology of strictly compact convergence. A fortiori, it is a Cauchy filter for vague convergence, and Prop. 14 of No.~9 shows that it is vaguely convergent to a measure \(\mu_{1}\) ; moreover, since V is dense in \(\mathcal{K}(\mathrm{X};\mathbf{C})\) , the hypothesis implies that \(\mu_{1} = \mu_{0}\) ; finally, since U is a Cauchy filter for the topology of strictly compact convergence, it also converges to \(\mu_{0}\) for this topology (GT, X, §1, No.~5, Prop. 5).

Q.E.D.

COROLLARY. --- If X is paracompact then the topologies induced on \(\mathcal{M}_{+}(X)\) by the vague topology and the topology of compact convergence coincide.

However, the topologies induced on \(\mathcal{M}_{+}(X)\) by the topology of compact convergence and the topology of strictly compact convergence may be different when X is not paracompact (Exer. 3).

